\documentclass{article}
\usepackage[T1]{fontenc}
\usepackage{lmodern}
\usepackage{graphicx}
\usepackage{amsmath,amssymb}
\usepackage[a4paper,margin=2cm]{geometry}
\usepackage[absolute,overlay]{textpos}
\setlength{\TPHorizModule}{1mm}
\setlength{\TPVertModule}{1mm}
\usepackage[indonesian]{babel}
\usepackage{hyperref}
\setlength{\emergencystretch}{2em}
% unit: Halaman 2

\begin{document}

% === Halaman 2 (llm) ===

\section*{Pendahuluan}

\begin{itemize}
    \item RSA merupakan algoritma kunci-publik yang paling terkenal dan paling banyak aplikasinya.
    \item Dibuat oleh tiga peneliti dari \textit{MIT (Massachussets Institute of Technology)}, yaitu Ronald \textcolor{red}{Rivest}, Adi \textcolor{red}{Shamir}, dan Leonard \textcolor{red}{Adleman}, pada tahun 1976.
    \item RSA = Rivest-Shamir-Adleman
    \item Keamanan algoritma RSA terletak pada sulitnya memfaktorkan bilangan bulat yang besar menjadi faktor-faktor prima.
\end{itemize}

\end{document}
